\chapter{Exercicis proposats}
\label{cap:exercicis}

\begin{enumerate}
\item Donats els conjunts següents%
\begin{equation*}
\begin{array}{l}
A=\left\{ x:x=2n\text{ \ i \ }n\in \mathbb{N}\right\} \\
B=\left\{ x:x=3n\text{ \ i \ }n\in \mathbb{N}\right\} \\
C=\left\{ x:x=4n-1\text{ \ i \ }n\in \mathbb{N}\right\} \\
D=\left\{ x:x=4n+1\text{ \ i \ }n\in \mathbb{N}\right\}%
\end{array}%
\end{equation*}%
Troba $A\cap B$, $C\cup D$, $C\cap D$, $\complement A$ i $A\cap (C\cup D)$.

Solució: $A\cap B$ és el conjunt de múltiples de 6, $C\cup
D=\left\{ x:x=2n+1\text{ \ i \ }n\in \mathbb{N}\right\} $, $C\cap
D=\emptyset $, $\complement A$ és el conjunt dels nombres imparells, i $%
A\cap (C\cup D)=\emptyset $.

\item Sigui $U$ l'univers dels conjunts $A,B$ i $C$. Simplifica les següents expressions:

\begin{enumerate}
\item[a)] $(A\cap B)\cup (A\cap \complement B)\cup (\complement A\cap B)\cup
(\complement A\cap \complement B)$

\item[b)] $\left[ (A\cup B)\cap \complement \left( B\cap \complement
A\right) \right] \cup \left[ (A\cap B)\cup \complement \left( B\cup
\complement A\right) \right] $

\item[c)] $(A\cap B)\cup (A\cap C)\cup \complement \left( \complement A\cup
\complement B\right) $

\item[d)] $(A\cap B)\cup (A\cap \complement B)\cup (\complement A\cap C)\cup
(A\cap C)$
\end{enumerate}

Solució: a) $U$, b) $A$, c) $A\cap (B\cup C)$ i d) $A\cup C$

\item Si $\#P\left( A\cup \mathcal{P}(A)\right) =8$, calcula $\#A$.

Solució: $\#A=1$

\item Una empresa ofereix places d'electricista, de mecànic i de fuster.
Sabem que 12 persones sol·liciten plaça
d'electricista, 12 de mecànic, 15 de fuster, 3 d'electricista i mecànic, 4 de mecànic i fuster, 5 d'electricista i fuster i, finalment, 1 sol·licita plaça de les tres coses. Calcula quanta
gent ha fet alguna sol·licitud.

Solució: 28 persones

\item En una reunió hi ha 25 persones que són mèdics, músics
o polítics. Hi ha 20 metges, 12 músics i 17 polítics. Hi ha 8
que són mèdics i músics, 12 que són mèdics i polítics i 11 que són músics i polítics. (a) Quants polítics són músics i metges alhora? (b) Quantes persones hi ha amb una sola
professió?

Solució: (a) 7 i (b) 8

\item D'un grup de 1000 persones, 950 porten rellotge, 750 porten paraigua,
800 porten corbata i 850 porten barret. Troba el nombre mínim de
persones que porten les quatre coses.

Solució: 350

\item Comprova si les següents col·leccions de
conjunts formen una partició de $\mathbb{N}$

\begin{enumerate}
\item[a)] $A=\left\{ 2n:n\in \mathbb{N}\right\} $ i $B=\left\{ 2n-1:n\in
\mathbb{N}\right\} $

\item[b)] $A$ és el conjunt dels nombres primers, $B=\left\{
2,4,6,8,9\right\} $ i $C=\left\{ x\in \mathbb{N}:x\geq 10\right\} $

\item[c)] $A=\left\{ 2n:n\in \mathbb{N}\right\} $, $B=\left\{ 2n-1:n\in
\mathbb{N}\right\} $ i $C=\left\{ 5n:n\in \mathbb{N}\right\} $
\end{enumerate}

Solució: a) Sí, b) No, c) No

\item En el conjunt $A=\{1,2,3,4,5,6\}$ es defineix la relació%
\begin{equation*}
\begin{array}{ccc}
x~R~y & \Longleftrightarrow & y\text{ és múltiple de }x%
\end{array}%
\end{equation*}%
(a) Escriu el graf de la relació i (b) estudia les seves propietats.

Solució: (a) $R=\left\{
\begin{array}{c}
(1,1),(1,2),(1,3),(1,4),(1,5),(1,6),(2,2), \\
(2,4),(2,6),(3,3),(3,6),(4,4),(5,5),(6,6)%
\end{array}%
\right\} $ (b) $R$ és reflexiva, antisimètrica i transitiva.

\item Quina relació binària sobre un conjunt és simètrica i
antisimètrica?

Solució: La relació d'igualtat.

\item De les següents relacions, quines són d'equivalència? I en
cas de ser-ho, quines són les seves classes?

\begin{enumerate}
\item[a)] «Tenir la mateixa altura» en el
conjunt dels alumnes d'una classe

\item[b)] «Ser equipol·lent» en el conjunt dels vectors fixos del pla

\item[c)] «Equidistar d'un punt fix
donat» en el conjunt dels punts del pla

\item[d)] «Estar alineats amb un punt fix
donat» en el conjunt de parells de punts del pla sense el
punt fix donat
\end{enumerate}

Solució: (a) És d'equivalència i els alumnes queden classificats
segons les seves altures (b) És d'equivalència i les classes són
els vectors lliures del pla (c) És d'equivalència i les classes són les circumferències de centre el punt fix donat (d) És
d'equivalència i les classes són les rectes que passen pel punt fix
donat sense contenir aquest punt.

\item En el conjunt dels nombres reals es defineix la relació%
\begin{equation*}
\begin{array}{ccc}
x\sim y & \Longleftrightarrow & \left\lfloor x\right\rfloor =\left\lfloor
y\right\rfloor%
\end{array}%
\end{equation*}%
on $\left\lfloor x\right\rfloor$ i $\left\lfloor y\right\rfloor$ indiquen la part sencera dels nombres reals $x$ i $y$. És una relació d'equivalència? Si ho és, quines són
les seves classes?

Solució: Observa que $x\sim y$ si i només si existeix un nombre
enter $n$ tal que $x,y\in [ n,n+1)$. És una relació d'equivalència i les classes són els intervals de la forma $[n,n+1)$ amb $%
n\in \mathbb{Z}$.

\item Esbrina si la relació «$x$ divideix $y$%
» és d'ordre en cadascun dels conjunts que s'indiquen
a continuació, i, en el cas que ho sigui, és parcial o total? Troba
els seus elements maximals i minimals.

\begin{enumerate}
\item[a)] En el conjunt dels nombres naturals $\mathbb{N}$

\item[b)] En el conjunt dels nombres enters $\mathbb{Z}$
\end{enumerate}

Solució: (a) És d'ordre parcial, $1$ és minimal i no hi ha
elements maximals. (b) No és d'ordre

\item En el conjunt dels nombres reals ordenat segons la relació d'ordre
usual $\leq $ es consideren els següents subconjunts (a) $A=[1,5]$, (b) $%
B=(-2,-1]$, (c) $C=(\pi ,2\pi )$, (d) $D=[2,+\infty )$, (e) $E=(-5,+\infty )$
i (f) $F=(-\infty ,0)$. Calcula, si existeixen, mínim, màxim, ínfim i suprem de cadascun d'ells.

Solució: (a) $\min A=\inf A=1$ i $\max A=\sup A=5$, (b) $\inf B=-2$ i $
\max B=\sup B=-1$, (c) $\inf C=\pi $ i $\sup C=2\pi $, (d) $\min D=\inf D=2$%
, (e) $\inf E=-5$, (f) $\sup F=0$

\item En el conjunt $A=\{2,3,5,6,15\}$ es defineix la relació%
\begin{equation*}
\begin{array}{ccc}
x~R~y & \Longleftrightarrow & y\text{ és múltiple de }x%
\end{array}%
\end{equation*}%
(a) És una relació d'ordre? És parcial o total? (b) Troba màxim, mínim, suprem i ínfim de $B=\{2,3,6,15\}$. (c) Troba màxim, mínim, suprem i ínfim de $A$. (d) Hi ha elements maximals i
minimals d'$A$?

Solució: (a) És d'ordre parcial (b) No existeixen $\max B$, $\min B$%
, $\sup B$, $\inf B$. (c) No existeixen $\max A$, $\min A$, $\sup A$, $\inf
A $. (d) $2,3,5$ són minimals i $6,15$ són maximals.

\item Es donen els conjunts $A=\{2,3,4,7,8\}$ i $B=\{1,2,3,4,5,7,9\}$. Sigui
$f:A~\longrightarrow ~B$ l'aplicació definida per
\begin{equation*}
f(x)=\left\{
\begin{array}{ll}
\frac{x}{2} & \text{si \ }x\text{ \ és parell} \\
x & \text{si \ }x\text{ \ és senar}%
\end{array}%
\right.
\end{equation*}%
Estudia quina classe d'aplicació s'obté.

Solució: És una aplicació injectiva

\item Donades les aplicacions $f:\mathbb{R}~\longrightarrow ~\mathbb{R}$
definida per $f(x)=e^{x}$ i $g:[0,+\infty )~\longrightarrow ~\mathbb{R}$
definida per $g(x)=\sqrt{x}$. (a) De quina classe d'aplicacions són $f$
i $g$? (b) Canvia els conjunts de sortida i d'arribada de l'aplicació $f$
perquè sigui bijectiva. (c) Canvia els conjunts de sortida i d'arribada
de l'aplicació $g$ perquè sigui bijectiva. (d) Calcula les
aplicacions inverses de les aplicacions dels apartats (b) i (c). (e) Calcula si
és possible $f\circ g$ i $g\circ f$.

Solució: (a) $f$ és injectiva però no exhaustiva i $g$ és
injectiva però no exhaustiva. (b) L'aplicació $f:\mathbb{R}%
~\longrightarrow ~(0,+\infty )$ definida per $f(x)=e^{x}$ és bijectiva.
(c) L'aplicació $g:[0,+\infty )~\longrightarrow ~[0,+\infty )$ definida
per $g(x)=\sqrt{x}$ és bijectiva. (d) $f^{-1}:(0,+\infty
)~\longrightarrow ~\mathbb{R}$ amb $f^{-1}(x)=\ln x$ i $g^{-1}:[0,+\infty
)~\longrightarrow ~[0,+\infty )$ amb $g^{-1}(x)=x^{2}$. (e) $f\circ
g:[0,+\infty )~\longrightarrow ~\mathbb{R}$ amb $(f\circ g)(x)=e^{\sqrt{x}}$
i $g\circ f:\mathbb{R}~\longrightarrow ~\mathbb{R}$ amb $(g\circ f)(x)=\sqrt{%
e^{x}}$ ja que $\func{Im}f=(0,+\infty )\subset [ 0,+\infty )$.

\item Donada l'aplicació $f:\mathbb{R}~\longrightarrow ~\mathbb{R}$
definida per $f(x)=2x+1$, calcula $f(A)$ i $f^{-1}(A)$ en els casos següents: (a) $A=[1,3]$, (b) $A=[-2,-1)$, (c) $A=[1,+\infty )$ i (d) $A=(-\infty
,-2)$.

Solució: (a) $f(A)=[3,7]$ i $f^{-1}(A)=[0,1]$, (b) $f(A)=[-3,-1)$ i $
f^{-1}(A)=[-3/2,-1)$, (c) $f(A)=[3,+\infty )$ i $f^{-1}(A)=[0,+\infty )$ i
(d) $f(A)=(-\infty ,-3)$ i $f^{-1}(A)=(-\infty ,-3/2)$.

\item Sigui $f:\mathbb{R}-\{1\}~\longrightarrow ~\mathbb{R}-\{3\}$ definida
per%
\begin{equation*}
f(x)=\frac{3x-1}{x-1}
\end{equation*}%
Demostra que és bijectiva i calcula l'aplicació inversa.

Solució: Cal demostrar que és injectiva i exhaustiva. L'aplicació
inversa és%
\begin{equation*}
f^{-1}(x)=\frac{x-1}{x-3}
\end{equation*}
\end{enumerate}
